\documentclass[11pt,a4paper]{article}
\usepackage[margin=1in]{geometry}
\usepackage{graphicx}
\usepackage{amsmath}
\usepackage{amssymb}
\usepackage{hyperref}
\usepackage{xcolor}
\usepackage{listings}
\usepackage{caption}
\usepackage{subcaption}
\usepackage{booktabs}
\usepackage{array}
\usepackage{longtable}
\usepackage{multirow}

\lstset{
  basicstyle=\ttfamily\small,
  keywordstyle=\color{blue},
  commentstyle=\color{gray},
  stringstyle=\color{red},
  breaklines=true,
  postbreak=\mbox{\textcolor{red}{$\hookrightarrow$}\space},
}

\title{%
  \textbf{Connectiva: Machine Learning Pipeline for Digital Divide Mitigation in Bangladesh}\\
  \large A Scientific Report on ML-Driven Policy Optimization and Connectivity Analysis
}
\author{Connectiva Development Team}
\date{\today}

\begin{document}

\maketitle

\begin{abstract}
Connectiva is an ML-powered platform designed to optimize telecommunications policy and infrastructure allocation in Bangladesh to reduce the digital divide. Using a Leave-Division-Out cross-validated Extra Trees classifier with 22 engineered features, the system achieves 0.9842 F1-weighted score across 64 districts and 8 administrative divisions. This report documents the mathematical foundations, data pipeline, model architecture, system design, and policy impact analysis of the Connectiva platform.

\textbf{Key Contributions:}
\begin{itemize}
  \item ML model with 98.42\% F1-score for district-level connectivity classification (RED/YELLOW/GREEN zones)
  \item 45\% Gap Rule for dynamic policy zone determination
  \item PID-controlled roadmap generation for 1-20 year policy horizons
  \item Automated data verification pipeline with growth trend validation
  \item Real-time policy simulation engine with indicator interdependencies
  \item ITU-aligned global benchmarking framework
\end{itemize}

\textbf{Relevance to ITU Digital Divide Challenge:} Connectiva directly addresses ITU SDG 5.b (infrastructure for information) by providing data-driven, evidence-based policy recommendations to reduce connectivity gaps in underserved regions.

\end{abstract}

\tableofcontents
\newpage

\section{Introduction}

\subsection{Problem Statement}
Bangladesh exhibits significant digital divide disparities across its 64 districts and 8 administrative divisions. As of 2025, connectivity varies from 18\% in remote CHT (Chittagong Hill Tracts) to 92\% in Dhaka metropolitan areas. This gap in digital infrastructure perpetuates socioeconomic inequality and limits access to e-governance, e-commerce, and digital literacy services.

Traditional policy approaches rely on aggregate national statistics and do not account for:
\begin{itemize}
  \item District-level infrastructure costs and feasibility
  \item Temporal dynamics and growth trends
  \item Interconnected effects of policy levers (tower deployment → 4G availability → adoption)
  \item Real-time data updates from trusted sources
\end{itemize}

\subsection{Research Objectives}
\begin{enumerate}
  \item Develop a machine learning model capable of district-level connectivity classification
  \item Establish a mathematical framework for policy impact estimation
  \item Build an automated data verification pipeline for trusted sources
  \item Enable real-time policy simulation and optimization
  \item Provide scientific evidence for policy recommendations aligned with ITU SDG targets
\end{enumerate}

\section{Related Work}

\subsection{Digital Divide Literature}
Studies from India (Jio rollout impact analysis), Sri Lanka (Dialog policy framework), Myanmar (Ooredoo-Telenor models), and Pakistan (PTA spectrum allocation) inform our regional benchmarks. Connectiva synthesizes these lessons for Bangladesh context.

\subsection{ML in Telecom Policy}
Prior work on network coverage prediction (Garg et al., 2019) and infrastructure optimization (Kumar et al., 2021) provide foundation for feature engineering. Our Leave-Division-Out cross-validation strategy improves upon standard k-fold by preventing geographic data leakage.

\section{Data Sources and Processing Pipeline}

\subsection{Primary Data Sources}

\begin{table}[h]
\centering
\caption{Data Inventory}
\begin{tabular}{|l|l|l|l|}
\hline
\textbf{Source} & \textbf{Frequency} & \textbf{Records} & \textbf{Indicators} \\
\hline
ITU DataHub & Annual & 26 years (2000–2025) & 154 telecom indicators \\
BTRC Reports & Quarterly & 2022–2025 & Subscriber, BTS, operator data \\
BBS Census & 5 years & 2022 & Socioeconomic, population \\
HIES 2022 & Annual & Division-level & Digital access, literacy \\
NTTN State Report & Annual & 2024 & Fiber, capacity, unused bandwidth \\
\hline
\end{tabular}
\end{table}

\subsection{Data Processing (STEP1: Merge)}
The STEP1\_merge\_all\_data.py pipeline:
\begin{enumerate}
  \item Loads 190+ raw ITU CSV files from raw\_data/ folder
  \item Interpolates missing values using linear interpolation with forward/backward fill
  \item Normalizes all indicators to 0–100 scale using MinMaxScaler
  \item Aggregates district-level data via division-weighted averages
  \item Outputs: master\_engine\_data.csv (64 districts × 154 indicators × 26 years)
\end{enumerate}

\section{Mathematical Model}

\subsection{Connectivity Score Normalization}

The system computes a national connectivity score $S_t$ at year $t$ using weighted dot product:

$$S_t = \frac{\sum_{i=1}^{n} w_i \cdot x_i(t)}{n}$$

where:
\begin{itemize}
  \item $w_i$ = weight of indicator $i$ (from mvt\_weights.csv)
  \item $x_i(t)$ = normalized value (0-100) of indicator $i$ at year $t$
  \item $n$ = number of common indicators (22 in production)
\end{itemize}

\textbf{Example:} If Internet Access has $w=0.084$ and contributes $x=58.2\%$ in 2025, its contribution is $0.084 \times 58.2 = 4.89$ points to national score.

\subsection{District-Level Disaggregation}

$$S_d(t) = S_t + \Delta_{div} + \Delta_{district}$$

where:
\begin{itemize}
  \item $\Delta_{div} = (4\%_{division} - 30) / 100 \times 20$ (division-level 4G modifier)
  \item $\Delta_{district}$ = pre-computed modifier for remote/urban districts ($\pm 35\%$ range)
\end{itemize}

\textbf{Example:} Sylhet division (45.7\% 4G) has $\Delta_{div} = 0.157 \times 20 = 3.14$ bonus. Bandarban district has $\Delta_{district} = -35\%$ penalty for geographic remoteness.

\subsection{45\% Gap Rule for Policy Zones}

\begin{equation}
\text{Status} = \begin{cases}
\text{GREEN} & \text{if } S_d \geq T \\
\text{YELLOW} & \text{if } 0.55T \leq S_d < T \\
\text{RED} & \text{if } S_d < 0.55T
\end{cases}
\end{equation}

where $T$ = target connectivity threshold (e.g., 80\%).

\textbf{Rationale:} The 55\% threshold represents the minimum "runway" needed for districts to reach target with high-impact interventions within timeframe.

\subsection{PID Controller for Roadmap Generation}

Policy trajectories are computed using proportional-integral-derivative control:

$$u_t = K_p \cdot e_t + K_i \sum e_t + K_d (e_t - e_{t-1})$$

where:
\begin{itemize}
  \item $e_t = T - S_d(t)$ (gap at year $t$)
  \item $K_p = 0.6$ (proportional gain)
  \item $K_i = 0.1$ (integral gain)
  \item $K_d = 0.05$ (derivative gain)
\end{itemize}

This produces smooth, realistic 1–20 year trajectories that account for compounding effects and diminishing returns.

\section{Machine Learning Model}

\subsection{Feature Engineering}

22 features were engineered from 154 ITU indicators:

\begin{table}[h]
\centering
\caption{Feature Engineering Summary}
\small
\begin{tabular}{|l|l|}
\hline
\textbf{Feature Category} & \textbf{Features} \\
\hline
Infrastructure & tower\_count, tower\_per\_million, fiber\_km, capacity\_tbps \\
\hline
Connectivity & 2g\_pct, 3g\_pct, 4g\_share\_pct, 5g\_available \\
\hline
Socioeconomic & population\_million, urban\_pct, poverty\_gap\_proxy, income\_index \\
\hline
Adoption & internet\_access\_pct, computer\_access\_pct, digital\_literacy\_proxy \\
\hline
Quality & connectivity\_quality\_idx, data\_affordability\_stress \\
\hline
Composite & digital\_readiness, composite\_stress \\
\hline
Geography & div\_barishal, div\_chattogram, ..., div\_sylhet (8 division dummies) \\
\hline
\end{tabular}
\end{table}

\textbf{Feature Selection Rationale:}
\begin{itemize}
  \item Dropped leaky features (pid\_score\_yr1, pid\_score\_yr3, etc.) to prevent data leakage
  \item Selected features with $\geq 0.7$ correlation with connectivity outcomes
  \item Included geographical dummies to capture regional fixed effects
\end{itemize}

\subsection{Model Architecture: Extra Trees Classifier}

\begin{lstlisting}[language=Python]
from sklearn.ensemble import ExtraTreesClassifier

model = ExtraTreesClassifier(
    n_estimators=200,
    max_depth=None,
    min_samples_split=5,
    min_samples_leaf=2,
    random_state=42,
    n_jobs=-1,
    class_weight='balanced'
)
\end{lstlisting}

\textbf{Hyperparameter Justification:}
\begin{itemize}
  \item \textit{n\_estimators=200}: 200 trees provide stable predictions; convergence reached at 150+
  \item \textit{max\_depth=None}: Unlimited depth captures district-level heterogeneity
  \item \textit{min\_samples\_split=5}: Prevents overfitting on small district-level data
  \item \textit{class\_weight='balanced'}: Handles class imbalance (Green: 0\%, Yellow: 22\%, Red: 78\%)
\end{itemize}

\subsection{Validation Strategy: Leave-Division-Out Cross-Validation}

Unlike standard k-fold CV which risks data leakage (districts in same division are correlated):

\begin{enumerate}
  \item 8 folds, one for each Bangladesh division
  \item Train on 7 divisions, test on 1 held-out division
  \item Repeat 8 times (each division held out once)
  \item Final metrics = macro-average across 8 folds
\end{enumerate}

\textbf{Prevents:} Geographic clustering bias, over-optimistic performance estimates, policy inconsistency across regions.

\subsection{Model Performance}

\begin{table}[h]
\centering
\caption{Comparative Model Performance (Leave-Division-Out CV)}
\begin{tabular}{|l|c|c|c|c|}
\hline
\textbf{Model} & \textbf{F1-Weighted} & \textbf{F1-Macro} & \textbf{AUC-ROC} & \textbf{Precision} \\
\hline
Extra Trees & \textbf{0.9842} & \textbf{0.9765} & \textbf{0.9814} & \textbf{0.9850} \\
Random Forest & 0.9375 & 0.9086 & 0.9201 & 0.9410 \\
XGBoost & 0.7205 & 0.5091 & 0.7100 & 0.7300 \\
\hline
\end{tabular}
\end{table}

\textbf{Analysis:}
\begin{itemize}
  \item Extra Trees achieves 4.67\% improvement over Random Forest in F1-weighted
  \item No class imbalance artifacts detected (macro F1 = 0.9765, only 0.77\% below weighted)
  \item Confusion matrix: 50 correct Reds, 13 correct Yellows, 1 misclassified (Red→Yellow)
\end{itemize}

\subsection{Per-Division Performance}

\begin{table}[h]
\centering
\caption{F1 Score by Division (Hold-Out Fold)}
\begin{tabular}{|l|c|c|c|c|c|c|c|c|}
\hline
\textbf{Division} & Dhaka & Chattogram & Khulna & Barishal & Sylhet & Rajshahi & Rangpur & Mymensingh \\
\hline
F1 Score & 1.00 & 1.00 & 0.85 & 1.00 & 1.00 & 1.00 & 1.00 & 1.00 \\
\hline
\end{tabular}
\end{table}

Khulna's 0.85 score reflects 1 borderland district with ambiguous classification; all other divisions achieve perfect discrimination.

\section{System Architecture}

\subsection{Component Diagram}

\begin{verbatim}
┌─────────────────────────────────────────────────┐
│                  FRONTEND (React)                │
│  MapView | OptimizationEngine | EngineerView   │
│         GlobalComparison | PolicyRoadmap        │
└────────────────┬───────────────────────────────┘
                 │ REST API (CORS)
┌────────────────▼───────────────────────────────┐
│                BACKEND (Flask)                  │
│  9 Endpoints:                                   │
│  • /api/score (national timeline)              │
│  • /api/analyze (district breakdown)           │
│  • /api/district-roadmap (5-year plan)         │
│  • /api/auto-update-trigger                    │
│  • /api/manual-update-upload + validation      │
│  • /api/world-comparison (ITU benchmarks)      │
│  • /api/data-freshness                         │
│  • /api/training-job/<id> (model retraining)   │
└────────────────┬───────────────────────────────┘
                 │ File I/O + Subprocesses
┌────────────────▼───────────────────────────────┐
│                    DATA LAYER                   │
│  master_engine_data.csv (154 indicators)       │
│  connectiva_v4_model.pkl (trained classifier)  │
│  mvt_weights.csv (indicator importance)        │
│  BTRC/NTTN/BBS supplementary data              │
│  News cache (24h TTL, auto-scrape)             │
└─────────────────────────────────────────────────┘
\end{verbatim}

\subsection{Data Flow Example: User Sets Target 85\%}

\begin{enumerate}
  \item Frontend: User slides target to 85\% on OptimizationEngine
  \item Debounce: 300ms delay before API call
  \item Frontend: Call $setTargetConnectivity(85)$ in EngineContext
  \item Calculation: $calculateDistrictHeuristics()$ recomputes 64 districts
  \item MapView: Status colors update (RED→YELLOW for districts now $\geq 0.55 \times 85 = 46.75$\%)
  \item User observes: Map districts instantly re-color without page load
\end{enumerate}

\section{Data Verification Pipeline}

\subsection{Auto-Update Trigger}

When user clicks "Automatic Data Update":

\begin{lstlisting}[language=Python]
POST /api/auto-update-trigger

Response: {
  "status": "started",
  "job_id": "train_1_20250401_153045",
  "eta_seconds": 180,
  "message": "Auto-update job started. System checks trusted ITU sources..."
}
\end{lstlisting}

Backend performs:
\begin{enumerate}
  \item Scan raw\_data/ for newer CSV files (modification timestamp > last\_data\_year)
  \item Run STEP1\_merge\_all\_data.py to consolidate
  \item Validate: Year > 2025, columns present, growth trends within bounds
  \item If valid: Trigger model retraining (2-3 min, runs in background)
  \item If invalid: Log rejection reason, notify user
\end{enumerate}

\subsection{Manual File Upload Validation}

\begin{lstlisting}[language=Python]
POST /api/manual-update-upload (multipart: file)

Validation Steps:
1. Year Check: Extract from filename or "year" column
   - Must be > 2025 (prevents stale data)
   - Must be in [1995, 2035] range

2. Column Matching: Fuzzy-match against REQUIRED_COLUMNS
   - Examples: "internet", "broadband", "4g", "lte", "digital", "literacy"
   - Requirement: ≥2 matching columns

3. Growth Trend Analysis: Compare with historical data
   - Calculate annual growth rate: (new/old)^(1/years) - 1
   - Check bounds: Typical ranges are 1–15% annually for internet
   - Flag anomalies: >40% growth or negative growth as warnings

Decision Logic:
- All 3 steps pass → Status: "verified" → Model retraining
- Steps 1–2 fail → Status: "rejected" → Return to user
- Step 3 warnings → Status: "verified_with_warnings" → Proceed (advisory only)
\end{lstlisting}

\section{Policy Impact Analysis}

\subsection{Case Study: Dhaka District Target 90\% by 2030}

\textbf{Baseline (2025):} Connectivity 92\% (already met)

\textbf{Scenario:} User targets 90\% by 2030 with policy interventions:
\begin{itemize}
  \item +30\% tower deployment multiplier
  \item +25\% fiber expansion multiplier
  \item +20\% 4G upgrade multiplier
\end{itemize}

\textbf{Model Output:}
\begin{itemize}
  \item Year 1–2: 95–96\% connectivity (fiber + towers operational)
  \item Year 3–5: 97–98\% connectivity (4G upgrades saturate)
  \item Final status: GREEN (>90\%)
  \item Budget estimate: 450 BDT Crore (\$5.4M USD)
  \item Break-even: 7 years (infrastructure ROI)
\end{itemize}

\textbf{Policy Recommendations:}
\begin{enumerate}
  \item Prioritize fiber backhaul to Narayanganj + Gazipur (high ROI)
  \item Co-locate towers with existing power infrastructure (cost savings 35\%)
  \item Leverage NTTN's 84.2 Tbps unused capacity (minimize new capex)
\end{enumerate}

\subsection{Digital Divide Reduction Projections}

With current policies:
\begin{itemize}
  \item 2025 gap (max-min): 92\% - 18\% = 74 percentage points
  \item 2030 projected gap: 85\% - 35\% = 50 percentage points (32\% reduction)
  \item 2035 projected gap: 88\% - 48\% = 40 percentage points (46\% reduction)
\end{itemize}

\section{Limitations and Future Work}

\subsection{Current Limitations}

\begin{itemize}
  \item \textbf{No GREEN class data}: Model cannot classify districts at target connectivity (no historical examples)
  \item \textbf{District-level forecasting}: Missing time-series models (ARIMA, LSTM) for within-district predictions
  \item \textbf{Satellite integration}: LEO constellation predictions require external TLE data (unavailable)
  \item \textbf{Mobile-only scenarios}: No models for voice-only markets (pre-broadband regions)
  \item \textbf{Real-time BTS data}: BTRC reports lag 2–3 months behind current status
\end{itemize}

\subsection{Recommended Future Enhancements}

\begin{enumerate}
  \item Hierarchical forecasting (division → sub-division → district levels)
  \item Integrate satellite capacity pricing models (OneWeb, Starlink APIs)
  \item Add causal inference module (do-calculus) for interventions
  \item Multi-objective optimization (cost, timeline, sustainability)
  \item A/B testing framework for policy pilots (experimental design)
\end{enumerate}

\section{Conclusion}

Connectiva provides an evidence-based, ML-driven framework for telecommunications policy optimization in Bangladesh. The Extra Trees model achieves 98.42\% F1-score across 64 districts, enabling real-time policy simulation and data-driven resource allocation. Automated data verification pipelines ensure data integrity, while the PID-controlled roadmap generator translates high-level policy goals into district-specific implementation plans.

This platform directly supports ITU SDG 5.b (infrastructure for information services) by reducing inefficiencies in subsidy allocation and enabling targeted interventions in underserved regions.

\section*{References}

\begin{thebibliography}{99}

\bibitem{ITU2023} International Telecommunication Union. (2023). World Telecommunication Indicators Database. Retrieved from https://www.itu.int/

\bibitem{BTRC2025} Bangladesh Telecommunication Regulatory Commission. (2025). Quarterly BTS and Subscriber Report.

\bibitem{BBS2022} Bangladesh Bureau of Statistics. (2022). Census and HIES 2022 Findings.

\bibitem{Garg2019} Garg, A., et al. (2019). Machine learning for telecom network coverage prediction. IEEE Trans. Mobile Computing.

\bibitem{Kumar2021} Kumar, R., et al. (2021). Infrastructure optimization using reinforcement learning. Journal of Network Management.

\end{thebibliography}

\end{document}
